\documentclass[11pt]{article}
\usepackage[margin=18mm,paperwidth=210mm,paperheight=297mm]{geometry}
\usepackage{booktabs}
\usepackage{tabularx}
\usepackage{array}
\usepackage[table]{xcolor}
\usepackage{caption}
\usepackage{microtype}
\usepackage[T1]{fontenc}
\usepackage{lmodern}

\definecolor{tablealt}{RGB}{245,247,250}
\definecolor{headergray}{RGB}{55,65,81}
\captionsetup{font=small,labelfont=bf,skip=8pt,justification=centering}

\newcolumntype{L}[1]{>{\raggedright\arraybackslash}p{#1}}
\newcolumntype{Y}{>{\raggedright\arraybackslash}X}
\newcolumntype{C}[1]{>{\centering\arraybackslash}p{#1}}

\begin{document}
\pagestyle{empty}
\vspace*{0pt}
\begin{center}
\setlength{\tabcolsep}{5pt}
\renewcommand{\arraystretch}{1.30}
\begin{minipage}{0.98\textwidth}
\captionof{table}{Functional MRI paradigm summary. All BOLD runs share TR 937 ms, TE 37 ms, flip angle 52$^{\circ}$, 2 mm isotropic voxels, and multiband acceleration factor 8. Durations are approximate (volumes $\times$ TR). Stimulation was monocular; the non-stimulated eye was occluded.}
\centering
\begin{tabularx}{\textwidth}{@{}L{2.6cm}Y C{2.2cm} C{2.2cm} C{2.4cm}@{}}
\toprule
\textcolor{headergray}{\textbf{BIDS task}} &
\textcolor{headergray}{\textbf{Description}} &
\textcolor{headergray}{\textbf{Runs/session}} &
\textcolor{headergray}{\textbf{Volumes/run}} &
\textcolor{headergray}{\textbf{Duration (min)}} \\
\midrule
{\ttfamily task-rest} & Fixation-only monocular resting-state & 1 & 320 & 5.0 \\
\rowcolor{tablealt} {\ttfamily task-movie} & Naturalistic movie segments with fixation overlay & 4 & 210 & 3.3 \\
{\ttfamily task-fmri} & Block-design grating/checkerboard stimulation & 4 & 226 & 3.5 \\
\rowcolor{tablealt} {\ttfamily task-control} & Short control EPI acquisitions & 3 & 20 & 0.3 \\
\bottomrule
\end{tabularx}
\end{minipage}
\end{center}
\end{document}
